\section{Domain Task 2: Phân tích Tỷ lệ lấp đầy theo thời gian}

Câu hỏi nghiệp vụ: ``Tỷ lệ lấp đầy phòng thay đổi như thế nào qua các
tháng, và tính mùa vụ ảnh hưởng thế nào đến từng quận?''

Mục đích: Xác định xu hướng mùa cao điểm, thấp điểm để từ đó có cái nhìn
chi tiết về hiệu suất kinh doanh qua thời gian của từng khu vực.

\subsection{Biểu đồ 1: Tỷ lệ lấp đầy theo tháng của các quận (Multi-line Chart)}

\subsubsection*{A. Thiết kế Idiom}

\begin{longtable}{@{} p{2.8cm} p{12cm} @{}}
\toprule
\textbf{Đặc điểm} & \textbf{Chi tiết} \\
\midrule
\endhead

\bottomrule
\endfoot

\textbf{Idiom} & Multi-line Chart (Biểu đồ đa đường) \\
\addlinespace[0.2cm]

\textbf{What} & Tháng (Month): O, Quận (Neighbourhood Group): C, Tỷ lệ lấp đầy (Avg. is Booked): Q \\
\addlinespace[0.2cm]

\textbf{Why} & produce (derive) $\rightarrow$ browse, lookup $\rightarrow$ summarize, compare
\begin{itemize}[nosep, leftmargin=*, topsep=0.1cm]
    \item Dẫn xuất (derive) dữ liệu tỷ lệ lấp đầy từ trạng thái đặt phòng theo ngày.
    \item Duyệt (browse) và tìm kiếm (lookup) xu hướng biến động lấp đầy theo tiến trình thời gian 12 tháng.
    \item Tóm tắt (summarize) và so sánh (compare) hiệu suất kinh doanh, tính mùa vụ giữa 5 quận để nhận diện sự khác biệt về nhu cầu thị trường.
\end{itemize} \\
\addlinespace[0.2cm]

\textbf{How} & \textbf{Encode:}
\begin{itemize}[nosep, leftmargin=*, topsep=0.1cm]
    \item Mark: Point, Line (Đường nối các điểm dữ liệu).
    \item Channel:
    \begin{itemize}[nosep, leftmargin=0.5cm]
        \item PosX: Trục thời gian tiến lên (Tháng -- Ordinal).
        \item PosY: Mức độ lấp đầy (Tỷ lệ \% -- Quantitative).
        \item Hue Color: Phân biệt 5 quận (Categorical).
    \end{itemize}
\end{itemize}

\vspace{0.2cm}
\textbf{Manipulate:}
\begin{itemize}[nosep, leftmargin=*, topsep=0.1cm]
    \item Hover + Highlight: Rê chuột vào một đường để highlight xu hướng của riêng quận đó, giúp khắc phục hiện tượng rối mắt khi các đường cắt nhau (line crossing).
    \item Selection: Chọn một quận trên Legend để quan sát riêng biệt dữ liệu của khu vực đó.
\end{itemize}

\vspace{0.2cm}
\textbf{Reduce:}
\begin{itemize}[nosep, leftmargin=*, topsep=0.1cm]
    \item Filter: Có thể lọc theo Năm (Year) để so sánh xu hướng lấp đầy giữa các giai đoạn khác nhau (ví dụ: trước và sau đại dịch).
\end{itemize} \\
\addlinespace[0.2cm]

\textbf{Scale} &
\begin{itemize}[nosep, leftmargin=*, topsep=0pt]
    \item Main key: 12 (Tháng trong năm).
    \item Categorical key: 5 (Nhóm các quận của NYC).
    \item Y-axis range: 0.0 $\rightarrow$ 0.6
\end{itemize} \\
\end{longtable}

\begin{figure}[H]
    \centering
    \safeincludegraphics[width=0.9\linewidth]{charts/domain_task2_chart1.png}
    \caption{Hình 2.1. Biểu đồ phân tích sự lấp đầy theo tháng của các quận}
    \label{fig:domain-task2-chart1}
\end{figure}

\subsubsection*{B. Đánh giá biểu đồ}

\textbf{1. Nguyên lý biểu đạt (Expressiveness):}

\begin{itemize}
    \item Thuộc tính: Tháng (O), Quận (C), Tỷ lệ lấp đầy (Q).
    \item Channel:
    \begin{itemize}
        \item PosX: tiến trình thời gian $\rightarrow$ dùng cho thuộc tính Ordinal (O).
        \item PosY: định lượng tỷ lệ phần trăm $\rightarrow$ dùng cho thuộc tính Quantitative (Q).
        \item Color (Hue): phân loại các quận $\rightarrow$ dùng cho thuộc tính Categorical (C).
    \end{itemize}
    \item Đánh giá mức độ quan trọng và phân bổ kênh theo Mackinlay:
    \begin{itemize}
        \item Ưu tiên 1 -- Vị trí trục tung (PosY): Sự biến động của Tỷ lệ lấp đầy (Q) là mục tiêu quan sát chính. Việc gán vào trục PosY tuân thủ đúng nguyên tắc: dùng kênh biểu đạt mạnh nhất (Position) cho biến định lượng cốt lõi nhất.
        \item Ưu tiên 2 -- Vị trí trục hoành (PosX): Thời gian/Tháng (O) là chiều quan trọng thứ hai để tạo thành chuỗi sự kiện. Dùng PosX tạo cảm giác tiến trình (sequential) là lựa chọn tối ưu theo quy luật nhận thức.
        \item Ưu tiên 3 -- Màu sắc (Color Hue): Phân nhóm các quận (C) đóng vai trò đối chiếu. Dùng Color Hue là lựa chọn chuẩn xác nhất cho dữ liệu phân loại (Categorical) khi kênh vị trí đã được dùng để vẽ line.
    \end{itemize}
    \item Kết luận: Áp dụng hoàn toàn chính xác các kênh biểu đạt. Line chart kết hợp với Hue Color là quy chuẩn tối ưu nhất để phân tích chuỗi thời gian có phân rã theo nhóm.
\end{itemize}

\textbf{2. Nguyên lý hiệu quả (Effectiveness):}

\begin{itemize}
    \item Độ chính xác (Accuracy): Trục PosY (vị trí điểm trên một thang đo chung) tuân theo Stevens' Psychophysical law với $n = 1.0$. Việc ước lượng độ lớn và so sánh khoảng cách giữa các điểm là cực kỳ chính xác, độ lỗi Log\_error là thấp nhất. Người xem dễ dàng nhận ra khoảng cách chênh lệch giữa các quận ở cùng một tháng.
    \item Khả năng phân biệt (Discriminability): Biểu đồ sử dụng 5 màu (Hue Color) cho 5 quận. Số lượng phân loại $< 7$ nên nhận thức màu của mắt hoàn toàn không bị ảnh hưởng, rất dễ phân biệt. Tuy nhiên, ở các giai đoạn tháng 8 đến tháng 10, các đường (lines) có hiện tượng tiệm cận và cắt nhau (line crossing/clutter), gây đôi chút khó khăn cho mắt, nhưng thao tác Highlight khi hover đã giúp giải quyết vấn đề này.
    \item Khả năng tách biệt (Separability): Kênh vị trí (Pos) và kênh màu sắc (Color) tách biệt hoàn toàn. Việc các điểm nằm ở vị trí cao/thấp không làm ảnh hưởng đến khả năng nhận diện màu sắc của quận đó.
\end{itemize}

\subsubsection*{C. Phân tích biểu đồ (Insight)}

\begin{itemize}
    \item Vì dữ liệu ban đầu được cào vào đầu tháng 11 năm 2025 nên ta thấy ở cả 3 quận tháng 11/2025 luôn là tháng có tỷ lệ đặt phòng cao nhất. Bên cạnh đó, khi nhìn vào giá trị đặt phòng trong tương lai thì ta cũng thấy sự vượt trội của tháng 11, người dân có vẻ chuẩn bị phòng kĩ càng cho tháng phục sinh vào cả năm sau.
    \item Manhattan (đường màu đỏ) thể hiện sự bứt phá mạnh mẽ vào cuối năm, đạt đỉnh cao nhất toàn thị trường vào tháng 11 (gần 60\%).
    \item Queens duy trì phong độ ổn định và dẫn đầu trong suốt giai đoạn giữa năm (tháng 5 -- tháng 10), trong khi Staten Island và Bronx luôn nằm ở nhóm dưới với tỷ lệ lấp đầy thấp.
\end{itemize}
